\documentclass[11pt,a4paper]{article}

%================================================
% PACKAGES
%================================================
\usepackage[utf8]{inputenc}
\usepackage[T1]{fontenc}
\usepackage[french]{babel}
\usepackage[a4paper,margin=2cm]{geometry}

\usepackage{amsmath,amssymb,amsthm,mathtools}
\usepackage{lmodern}
\usepackage{xcolor}
\usepackage{hyperref}
\usepackage{fancyhdr}
\usepackage{lastpage}
\usepackage{enumitem}
\usepackage{tcolorbox}
\usepackage{tikz}
\usepackage{array}
\usepackage{booktabs}
\usepackage{multicol}
\usepackage{titlesec}
\usepackage{siunitx}
\usepackage{physics}
\usepackage{chemfig}

%================================================
% MISE EN PAGE
%================================================
\pagestyle{fancy}
\fancyhf{}
\fancyhead[L]{Physique-Chimie -- Spécialité}
\fancyhead[R]{Exercices : cohésion de la matière et spectroscopie IR}
\fancyfoot[C]{\thepage\ /\pageref{LastPage}}

\hypersetup{
    colorlinks=true,
    linkcolor=blue!60!black,
    urlcolor=blue!60!black,
    pdftitle={Exercices corrigés - Cohésion de la matière et spectroscopie IR},
    pdfauthor={},
    pdfsubject={Physique-Chimie Spécialité}
}

\setlength{\parindent}{0pt}
\setlength{\parskip}{0.45em}

\titleformat{\section}{\Large\bfseries\color{blue!60!black}}{\thesection.}{0.6em}{}
\titleformat{\subsection}{\large\bfseries\color{blue!45!black}}{\thesubsection.}{0.6em}{}
\titleformat{\subsubsection}{\normalsize\bfseries\color{blue!30!black}}{\thesubsubsection.}{0.6em}{}

%================================================
% BOITES
%================================================
\newtcolorbox{objectifbox}{
    colback=blue!5,
    colframe=blue!60!black,
    title=Objectif,
    fonttitle=\bfseries
}

\newtcolorbox{enoncebox}{
    colback=gray!4,
    colframe=black!70,
    title=Énoncé,
    fonttitle=\bfseries
}

\newtcolorbox{corrigebox}{
    colback=green!5,
    colframe=green!50!black,
    title=Corrigé,
    fonttitle=\bfseries
}

\newtcolorbox{ideebox}{
    colback=orange!8,
    colframe=orange!70!black,
    title=Idée essentielle,
    fonttitle=\bfseries
}

\newtcolorbox{attentionbox}{
    colback=red!5,
    colframe=red!65!black,
    title=Point de vigilance,
    fonttitle=\bfseries
}

\newtcolorbox{methodebox}{
    colback=purple!5,
    colframe=purple!60!black,
    title=Méthode,
    fonttitle=\bfseries
}

\newtcolorbox{bilanbox}{
    colback=cyan!6,
    colframe=cyan!55!black,
    title=Bilan,
    fonttitle=\bfseries
}

%================================================
% COMMANDES
%================================================
\newcommand{\chapitreA}{cohésion de la matière}
\newcommand{\chapitreB}{spectroscopie infrarouge}

%================================================
% DOCUMENT
%================================================
\begin{document}

\begin{center}
    {\LARGE \textbf{Feuille d'exercices approfondis}}\\[0.3em]
    {\Large \textbf{Cohésion de la matière et spectroscopie IR}}\\[0.4em]
    {\normalsize Avec corrigé détaillé sur le même document}\\[0.7em]
    \rule{0.9\linewidth}{0.6pt}
\end{center}

\tableofcontents
\vspace{1em}

%================================================
\section{Introduction}
%================================================

\begin{objectifbox}
Cette feuille ne vise pas seulement à reconnaître mécaniquement une liaison, une molécule polaire ou une bande IR.  
Elle cherche surtout à faire comprendre :
\begin{itemize}
    \item \textbf{pourquoi} certaines interactions existent ;
    \item \textbf{comment} elles expliquent les propriétés observables ;
    \item \textbf{pourquoi} une interprétation IR a du sens physiquement ;
    \item \textbf{comment relier} structure, interactions, propriétés et analyse spectrale.
\end{itemize}
\end{objectifbox}

%================================================
\section{Questions de cours réfléchies}
%================================================

%------------------------------------------------
\subsection{Question 1 -- Liaison et interaction}
%------------------------------------------------

\begin{enoncebox}
On considère les affirmations suivantes :
\begin{enumerate}[label=\alph*)]
    \item Une liaison covalente et une interaction de Van der Waals sont deux phénomènes de même nature, seule leur intensité change.
    \item Dans une molécule, la liaison covalente intervient à l'intérieur de l'espèce, alors que les interactions de Van der Waals agissent entre espèces distinctes.
    \item Une liaison hydrogène est une liaison covalente faible.
    \item Une interaction intermoléculaire peut modifier des propriétés macroscopiques d'une substance.
\end{enumerate}

Pour chaque affirmation, indiquer si elle est vraie ou fausse, puis justifier soigneusement.
\end{enoncebox}

\begin{corrigebox}
\begin{enumerate}[label=\alph*)]
    \item \textbf{Faux.}  
    Une liaison covalente est une liaison chimique intramoléculaire fondée sur un partage d'électrons.  
    Une interaction de Van der Waals est une interaction intermoléculaire faible entre molécules distinctes.  
    Ce n'est donc pas seulement une question d'intensité : la nature physique et le rôle structural ne sont pas les mêmes.
    
    \item \textbf{Vrai.}  
    C'est précisément une distinction centrale du chapitre :
    \begin{itemize}
        \item la liaison covalente maintient ensemble les atomes dans une même molécule ;
        \item les interactions de Van der Waals interviennent entre deux molécules distinctes.
    \end{itemize}
    
    \item \textbf{Faux.}  
    La liaison hydrogène n'est pas une liaison covalente.  
    C'est une interaction intermoléculaire attractive particulière, plus forte que beaucoup d'interactions de Van der Waals, mais plus faible qu'une liaison covalente.
    
    \item \textbf{Vrai.}  
    Les interactions intermoléculaires influencent de nombreuses propriétés :
    \begin{itemize}
        \item température d'ébullition ;
        \item température de fusion ;
        \item viscosité ;
        \item solubilité ;
        \item miscibilité.
    \end{itemize}
\end{enumerate}

\begin{ideebox}
La vraie frontière conceptuelle n'est pas seulement ``fort / faible'', mais surtout :
\[
\text{dans l'entité} \quad \text{ou} \quad \text{entre les entités.}
\]
\end{ideebox}
\end{corrigebox}

%------------------------------------------------
\subsection{Question 2 -- Polarité et géométrie}
%------------------------------------------------

\begin{enoncebox}
Expliquer pourquoi la phrase suivante est incomplète :

\begin{center}
\textit{``Une molécule est polaire dès qu'elle contient une liaison polaire.''}
\end{center}

On demande :
\begin{enumerate}[label=\arabic*)]
    \item d'expliquer précisément ce qui manque dans cette phrase ;
    \item de donner un exemple de molécule ayant des liaisons polaires mais étant globalement apolaire ;
    \item de donner un exemple de molécule polaire ;
    \item de formuler une version correcte de la phrase.
\end{enumerate}
\end{enoncebox}

\begin{corrigebox}
\begin{enumerate}[label=\arabic*)]
    \item Cette phrase est incomplète car elle oublie le rôle de la \textbf{géométrie} de la molécule.  
    Une liaison polaire crée un déséquilibre local de charge, mais les effets de plusieurs liaisons peuvent se compenser si la disposition spatiale est symétrique.
    
    \item L'exemple classique est le dioxyde de carbone \(\mathrm{CO_2}\).  
    Chaque liaison C=O est polaire, car l'oxygène est plus électronégatif que le carbone.  
    Mais la molécule est linéaire et symétrique ; les effets des deux liaisons se compensent.  
    La molécule est donc globalement apolaire.
    
    \item L'eau \(\mathrm{H_2O}\) est un exemple de molécule polaire.  
    Les liaisons O--H sont polaires et la molécule est coudée, donc les effets ne se compensent pas.
    
    \item Version correcte :
    
    \begin{center}
    \textit{Une molécule peut être polaire si elle contient des liaisons polaires et si la géométrie de la molécule ne compense pas leurs effets.}
    \end{center}
\end{enumerate}

\begin{attentionbox}
Une liaison polaire ne suffit jamais à elle seule pour conclure sur la polarité globale d'une molécule.
\end{attentionbox}
\end{corrigebox}

%------------------------------------------------
\subsection{Question 3 -- Le sens physique de la dissolution}
%------------------------------------------------

\begin{enoncebox}
On lit souvent : \textit{« le sel se dissout dans l'eau »}.  

Cette phrase peut donner l'impression que le sel ``disparaît''.

\begin{enumerate}[label=\arabic*)]
    \item Expliquer pourquoi cette idée est fausse.
    \item Décrire microscopiquement ce qui se passe lors de la dissolution d'un solide ionique dans l'eau.
    \item Expliquer pourquoi la polarité de l'eau joue un rôle décisif.
\end{enumerate}
\end{enoncebox}

\begin{corrigebox}
\begin{enumerate}[label=\arabic*)]
    \item Le sel ne disparaît pas.  
    Il cesse d'être visible comme cristal macroscopique, mais ses constituants sont toujours présents dans la solution sous forme d'ions dispersés.
    
    \item Dans un cristal ionique comme \(\mathrm{NaCl}\), les ions \(\mathrm{Na^+}\) et \(\mathrm{Cl^-}\) sont organisés dans un réseau.  
    Lorsqu'on met le cristal dans l'eau, les molécules d'eau entourent progressivement les ions, ce qui permet de les extraire du réseau cristallin.  
    On obtient alors des ions hydratés dispersés dans tout le solvant.
    
    \item La polarité de l'eau est essentielle car :
    \begin{itemize}
        \item le côté partiellement négatif de l'eau attire les cations ;
        \item le côté partiellement positif attire les anions.
    \end{itemize}
    Ces interactions eau--ions stabilisent les ions en solution et rendent la dissolution possible.
\end{enumerate}

\begin{ideebox}
Se dissoudre, ce n'est pas s'anéantir : c'est passer d'un état ordonné collectif à un état dispersé, mais stabilisé par de nouvelles interactions.
\end{ideebox}
\end{corrigebox}

%================================================
\section{Exercices de compréhension fine}
%================================================

%------------------------------------------------
\subsection{Exercice 1 -- Comparer trois substances}
%------------------------------------------------

\begin{enoncebox}
On considère trois substances :
\begin{itemize}
    \item l'eau \(\mathrm{H_2O}\),
    \item le diiode \(\mathrm{I_2}\),
    \item le chlorure de sodium \(\mathrm{NaCl}\).
\end{itemize}

Pour chacune d'elles :
\begin{enumerate}[label=\arabic*)]
    \item préciser la nature des entités qui la constituent ;
    \item indiquer le type principal de cohésion ;
    \item dire si l'eau liquide peut dissoudre cette substance, et justifier qualitativement ;
    \item expliquer pourquoi ces trois cas sont très différents du point de vue microscopique.
\end{enumerate}
\end{enoncebox}

\begin{corrigebox}
\textbf{1) Nature des entités}

\begin{itemize}
    \item \(\mathrm{H_2O}\) : molécules d'eau.
    \item \(\mathrm{I_2}\) : molécules de diiode.
    \item \(\mathrm{NaCl}\) : ions \(\mathrm{Na^+}\) et \(\mathrm{Cl^-}\) organisés en réseau cristallin.
\end{itemize}

\textbf{2) Type principal de cohésion}

\begin{itemize}
    \item \(\mathrm{H_2O}\) : interactions intermoléculaires, notamment liaisons hydrogène.
    \item \(\mathrm{I_2}\) : interactions de Van der Waals entre molécules apolaires.
    \item \(\mathrm{NaCl}\) : interaction électrostatique entre ions de charges opposées.
\end{itemize}

\textbf{3) Dissolution dans l'eau}

\begin{itemize}
    \item \(\mathrm{H_2O}\) se mélange évidemment avec lui-même.
    \item \(\mathrm{I_2}\) est peu soluble dans l'eau, car il s'agit d'une molécule apolaire, alors que l'eau est polaire. Les nouvelles interactions créées sont peu favorables.
    \item \(\mathrm{NaCl}\) se dissout bien dans l'eau, car l'eau peut hydrater les ions.
\end{itemize}

\textbf{4) Pourquoi ces cas sont-ils très différents ?}

Ils correspondent à trois logiques microscopiques distinctes :
\begin{itemize}
    \item dans l'eau, ce sont des molécules polaires capables de faire des liaisons hydrogène ;
    \item dans le diiode, ce sont des molécules apolaires attirées faiblement ;
    \item dans le chlorure de sodium, il n'y a pas de molécules isolées mais un réseau ionique.
\end{itemize}

\begin{bilanbox}
Cet exercice oblige à distinguer soigneusement :
\begin{itemize}
    \item molécule polaire,
    \item molécule apolaire,
    \item solide ionique.
\end{itemize}
C'est une distinction beaucoup plus profonde que la simple récitation d'une définition.
\end{bilanbox}
\end{corrigebox}

%------------------------------------------------
\subsection{Exercice 2 -- Température d'ébullition et interactions}
%------------------------------------------------

\begin{enoncebox}
On compare qualitativement les températures d'ébullition de quatre espèces :
\[
\mathrm{CH_4},\quad \mathrm{H_2O},\quad \mathrm{HCl},\quad \mathrm{NH_3}
\]

On rappelle seulement les idées suivantes :
\begin{itemize}
    \item \(\mathrm{CH_4}\) est une molécule apolaire ;
    \item \(\mathrm{HCl}\) est une molécule polaire ;
    \item \(\mathrm{NH_3}\) et \(\mathrm{H_2O}\) peuvent former des liaisons hydrogène, l'eau davantage.
\end{itemize}

\begin{enumerate}[label=\arabic*)]
    \item Classer qualitativement ces espèces selon la force croissante des interactions intermoléculaires.
    \item En déduire un classement plausible de leurs températures d'ébullition.
    \item Expliquer pourquoi une petite molécule comme l'eau peut avoir une température d'ébullition relativement élevée.
\end{enumerate}
\end{enoncebox}

\begin{corrigebox}
\textbf{1) Force croissante des interactions}

\begin{itemize}
    \item \(\mathrm{CH_4}\) : molécule apolaire \(\Rightarrow\) essentiellement interactions de Van der Waals faibles.
    \item \(\mathrm{HCl}\) : molécule polaire \(\Rightarrow\) interactions dipôle-dipôle plus fortes.
    \item \(\mathrm{NH_3}\) : possibilité de liaisons hydrogène.
    \item \(\mathrm{H_2O}\) : liaisons hydrogène nombreuses et très marquées.
\end{itemize}

Un classement raisonnable est donc :
\[
\mathrm{CH_4} < \mathrm{HCl} < \mathrm{NH_3} < \mathrm{H_2O}
\]

\textbf{2) Températures d'ébullition}

Plus les interactions intermoléculaires sont fortes, plus il faut fournir d'énergie pour séparer les molécules dans l'état gazeux.  
On s'attend donc à :
\[
T_{\text{éb}}(\mathrm{CH_4}) < T_{\text{éb}}(\mathrm{HCl}) < T_{\text{éb}}(\mathrm{NH_3}) < T_{\text{éb}}(\mathrm{H_2O})
\]

\textbf{3) Pourquoi l'eau bout-elle ``haut'' ?}

L'eau est petite, mais ce n'est pas la taille seule qui gouverne ici.  
Sa forte polarité et surtout sa capacité à établir de nombreuses liaisons hydrogène renforcent fortement la cohésion du liquide.  
Il faut donc fournir beaucoup d'énergie pour rompre cette organisation intermoléculaire.

\begin{ideebox}
Cet exercice montre qu'une propriété macroscopique comme la température d'ébullition n'est intelligible qu'en revenant aux interactions microscopiques.
\end{ideebox}
\end{corrigebox}

%------------------------------------------------
\subsection{Exercice 3 -- Polarité ou non ?}
%------------------------------------------------

\begin{enoncebox}
On considère les molécules suivantes :
\[
\mathrm{H_2O},\quad \mathrm{CO_2},\quad \mathrm{CH_4},\quad \mathrm{NH_3}
\]

On admet les géométries suivantes :
\begin{itemize}
    \item \(\mathrm{H_2O}\) : coudée,
    \item \(\mathrm{CO_2}\) : linéaire,
    \item \(\mathrm{CH_4}\) : tétraédrique symétrique,
    \item \(\mathrm{NH_3}\) : pyramidale.
\end{itemize}

\begin{enumerate}[label=\arabic*)]
    \item Dire quelles molécules sont polaires.
    \item Pour chaque molécule, expliquer le raisonnement au lieu de simplement donner la réponse.
    \item Expliquer pourquoi \(\mathrm{CO_2}\) et \(\mathrm{H_2O}\), qui contiennent tous deux des liaisons avec l'oxygène, n'ont pas le même comportement global.
\end{enumerate}
\end{enoncebox}

\begin{corrigebox}
\textbf{1) Molécules polaires}

\begin{itemize}
    \item \(\mathrm{H_2O}\) : polaire.
    \item \(\mathrm{CO_2}\) : apolaire.
    \item \(\mathrm{CH_4}\) : apolaire.
    \item \(\mathrm{NH_3}\) : polaire.
\end{itemize}

\textbf{2) Raisonnement détaillé}

\begin{itemize}
    \item \(\mathrm{H_2O}\) : les liaisons O--H sont polaires, et la géométrie coudée empêche la compensation \(\Rightarrow\) molécule polaire.
    \item \(\mathrm{CO_2}\) : les liaisons C=O sont polaires, mais la géométrie linéaire place les deux effets en opposition parfaite \(\Rightarrow\) compensation \(\Rightarrow\) molécule apolaire.
    \item \(\mathrm{CH_4}\) : géométrie tétraédrique symétrique ; l'ensemble est globalement équilibré \(\Rightarrow\) molécule apolaire.
    \item \(\mathrm{NH_3}\) : les liaisons N--H sont polaires, et la géométrie pyramidale n'assure pas de compensation \(\Rightarrow\) molécule polaire.
\end{itemize}

\textbf{3) Comparaison entre \(\mathrm{CO_2}\) et \(\mathrm{H_2O}\)}

La différence vient de la géométrie :
\begin{itemize}
    \item \(\mathrm{CO_2}\) est linéaire et symétrique ;
    \item \(\mathrm{H_2O}\) est coudée.
\end{itemize}
Ainsi, la polarité des liaisons n'a pas le même effet global.

\begin{attentionbox}
Deux molécules peuvent contenir un même type de liaison polaire tout en ayant une polarité globale différente.  
La géométrie reste indispensable.
\end{attentionbox}
\end{corrigebox}

%------------------------------------------------
\subsection{Exercice 4 -- Solubilité : raisonnement plutôt que recette}
%------------------------------------------------

\begin{enoncebox}
On souhaite prévoir qualitativement la solubilité dans l'eau des espèces suivantes :
\[
\mathrm{NaCl},\quad \mathrm{I_2},\quad \text{éthanol},\quad \text{huile}
\]

On rappelle que l'éthanol possède un groupe O--H.

\begin{enumerate}[label=\arabic*)]
    \item Dire, pour chacune, si elle est très soluble, peu soluble ou non miscible / peu miscible avec l'eau.
    \item Justifier à chaque fois en termes d'interactions microscopiques.
    \item Expliquer pourquoi la règle ``le semblable dissout le semblable'' est utile, mais non suffisante si elle est récitéée sans réflexion.
\end{enumerate}
\end{enoncebox}

\begin{corrigebox}
\textbf{1) Prévision qualitative}

\begin{itemize}
    \item \(\mathrm{NaCl}\) : très soluble.
    \item \(\mathrm{I_2}\) : peu soluble.
    \item Éthanol : miscible ou très soluble.
    \item Huile : non miscible ou très peu miscible.
\end{itemize}

\textbf{2) Justifications microscopiques}

\begin{itemize}
    \item \(\mathrm{NaCl}\) : l'eau, molécule polaire, hydrate les ions \(\mathrm{Na^+}\) et \(\mathrm{Cl^-}\).
    \item \(\mathrm{I_2}\) : molécule apolaire ; les interactions avec l'eau sont peu favorables.
    \item Éthanol : le groupe O--H rend possibles des interactions favorables avec l'eau, notamment des liaisons hydrogène.
    \item Huile : constituée en grande partie d'espèces apolaires ; les interactions avec l'eau sont défavorables.
\end{itemize}

\textbf{3) Limite de la règle}

La règle ``le semblable dissout le semblable'' est une bonne première idée directrice, mais elle doit être traduite en termes d'interactions :
\begin{itemize}
    \item quelles interactions faut-il rompre ?
    \item quelles interactions nouvelles se forment ?
    \item ces nouvelles interactions sont-elles assez stabilisantes ?
\end{itemize}
Sans cela, la règle devient un slogan sans contenu physique.

\begin{ideebox}
Une règle utile n'est vraiment comprise que lorsqu'on sait la relier à un mécanisme.
\end{ideebox}
\end{corrigebox}

%================================================
\section{Exercices sur la spectroscopie IR}
%================================================

%------------------------------------------------
\subsection{Exercice 5 -- Lire un spectre sans image}
%------------------------------------------------

\begin{enoncebox}
On décrit trois spectres IR fictifs :

\begin{itemize}
    \item \textbf{Spectre A} : une bande large vers \(3300\ \si{cm^{-1}}\), pas de bande forte vers \(1700\ \si{cm^{-1}}\).
    \item \textbf{Spectre B} : une bande forte vers \(1710\ \si{cm^{-1}}\), pas de bande large entre \(3200\) et \(3600\ \si{cm^{-1}}\).
    \item \textbf{Spectre C} : une bande forte vers \(1700\ \si{cm^{-1}}\) et une très large bande entre \(2500\) et \(3200\ \si{cm^{-1}}\).
\end{itemize}

Pour chacun :
\begin{enumerate}[label=\arabic*)]
    \item identifier la ou les fonctions probables ;
    \item expliquer le raisonnement ;
    \item dire pourquoi il serait excessif de prétendre identifier la molécule exacte avec ces seules données.
\end{enumerate}
\end{enoncebox}

\begin{corrigebox}
\textbf{Spectre A}

La bande large vers \(3300\ \si{cm^{-1}}\) évoque un groupe O--H d'alcool.  
L'absence de bande forte vers \(1700\ \si{cm^{-1}}\) indique qu'il n'y a probablement pas de groupe carbonyle.  
La fonction probable est donc \textbf{alcool}.

\textbf{Spectre B}

La bande forte vers \(1710\ \si{cm^{-1}}\) indique un groupe C=O.  
Comme il n'y a pas de grande bande O--H, on ne pense pas à un alcool ni à un acide carboxylique.  
On peut envisager une \textbf{cétone} ou un \textbf{aldéhyde}.

\textbf{Spectre C}

La bande C=O vers \(1700\ \si{cm^{-1}}\) signale un carbonyle.  
La très large bande entre \(2500\) et \(3200\ \si{cm^{-1}}\) est caractéristique d'un O--H acide.  
On conclut à un \textbf{acide carboxylique}.

\textbf{Pourquoi ne peut-on pas identifier la molécule exacte ?}

Un spectre IR permet surtout d'identifier des \textbf{groupes caractéristiques}.  
Plusieurs molécules différentes peuvent posséder la même fonction.  
Par exemple, de nombreux alcools ont un spectre présentant une bande O--H.  
On peut donc identifier une \textbf{famille} ou une \textbf{fonction probable}, mais pas nécessairement la structure exacte complète.

\begin{attentionbox}
Un spectre IR est un outil d'identification partielle, pas une carte d'identité unique dans tous les cas.
\end{attentionbox}
\end{corrigebox}

%------------------------------------------------
\subsection{Exercice 6 -- Pourquoi une bande est-elle large ?}
%------------------------------------------------

\begin{enoncebox}
Dans beaucoup d'exercices, on dit :
\begin{center}
\textit{« La bande O--H est large. »}
\end{center}

\begin{enumerate}[label=\arabic*)]
    \item Expliquer pourquoi cette largeur constitue déjà une information physique.
    \item Relier cette largeur à la présence possible de liaisons hydrogène.
    \item Expliquer pourquoi cette observation relie le chapitre de spectroscopie IR à celui de la cohésion de la matière.
\end{enumerate}
\end{enoncebox}

\begin{corrigebox}
\begin{enumerate}[label=\arabic*)]
    \item La largeur d'une bande signifie que toutes les liaisons O--H observées ne vibrent pas exactement dans le même environnement.  
    Cela montre qu'il existe une diversité de situations microscopiques.
    
    \item Lorsqu'un groupe O--H participe à des liaisons hydrogène, son environnement électronique est modifié.  
    Les vibrations possibles sont alors affectées, ce qui élargit souvent la zone d'absorption.
    
    \item Cette observation relie directement les deux chapitres :
    \begin{itemize}
        \item en \chapitreB{}, on observe une bande caractéristique ;
        \item en \chapitreA{}, on comprend que cette bande dépend des interactions intermoléculaires, ici les liaisons hydrogène.
    \end{itemize}
    Le spectre IR contient donc parfois une information indirecte sur la cohésion intermoléculaire.
\end{enumerate}

\begin{ideebox}
Le spectre IR ne révèle pas seulement qu'une liaison existe ; il peut aussi renseigner sur son environnement.
\end{ideebox}
\end{corrigebox}

%------------------------------------------------
\subsection{Exercice 7 -- Analyse critique d'une conclusion trop rapide}
%------------------------------------------------

\begin{enoncebox}
Un élève affirme :

\begin{center}
\textit{« J'observe une bande vers \(1700\ \si{cm^{-1}}\), donc la molécule est forcément un acide carboxylique. »}
\end{center}

\begin{enumerate}[label=\arabic*)]
    \item Expliquer pourquoi ce raisonnement est faux.
    \item Indiquer ce que signifie réellement une bande vers \(1700\ \si{cm^{-1}}\).
    \item Donner des exemples de fonctions différentes pouvant contenir ce type de liaison.
    \item Dire quelle information supplémentaire il faudrait rechercher pour conclure à un acide carboxylique.
\end{enumerate}
\end{enoncebox}

\begin{corrigebox}
\begin{enumerate}[label=\arabic*)]
    \item Le raisonnement est faux car une bande vers \(1700\ \si{cm^{-1}}\) ne suffit pas à elle seule à identifier une fonction unique.
    
    \item Cette bande indique surtout la présence probable d'un \textbf{groupe carbonyle} C=O.
    
    \item On trouve un groupe carbonyle dans :
    \begin{itemize}
        \item les aldéhydes ;
        \item les cétones ;
        \item les acides carboxyliques ;
        \item les esters.
    \end{itemize}
    
    \item Pour conclure à un acide carboxylique, il faut rechercher en plus une très large bande O--H acide entre \(2500\) et \(3200\ \si{cm^{-1}}\).
\end{enumerate}

\begin{attentionbox}
Une bande isolée conduit rarement à une conclusion certaine.  
L'analyse IR est un croisement d'indices.
\end{attentionbox}
\end{corrigebox}

%================================================
\section{Problèmes approfondis}
%================================================

%------------------------------------------------
\subsection{Problème 1 -- Pourquoi l'huile et l'eau se séparent-elles ?}
%------------------------------------------------

\begin{enoncebox}
On mélange vigoureusement de l'eau et de l'huile dans un récipient transparent.  
Pendant quelques instants, on observe une dispersion de fines gouttelettes, puis les deux phases se séparent à nouveau.

On demande une explication qualitative et cohérente de ce phénomène.

\begin{enumerate}[label=\arabic*)]
    \item Expliquer pourquoi l'agitation mécanique peut temporairement disperser l'une dans l'autre.
    \item Expliquer pourquoi cette dispersion n'est pas stable.
    \item Justifier pourquoi l'eau préfère interagir avec elle-même plutôt qu'avec l'huile.
    \item Expliquer en quoi cette expérience illustre les notions de polarité, de cohésion et de miscibilité.
\end{enumerate}
\end{enoncebox}

\begin{corrigebox}
\textbf{a) Effet de l'agitation}

L'agitation apporte de l'énergie mécanique et fragmente un liquide en fines gouttelettes dispersées dans l'autre.  
On peut donc obtenir temporairement une apparence de mélange.

\textbf{b) Pourquoi cette dispersion n'est-elle pas stable ?}

La dispersion n'est pas favorable énergétiquement à long terme, car les interactions eau--huile sont moins favorables que les interactions eau--eau et huile--huile.  
Le système évolue donc spontanément vers une séparation en deux phases.

\textbf{c) Pourquoi l'eau préfère-t-elle interagir avec elle-même ?}

L'eau est polaire et peut former des liaisons hydrogène entre ses molécules.  
L'huile est globalement apolaire.  
Les interactions entre eau et huile sont donc nettement moins favorables que les interactions internes à chaque liquide.

\textbf{d) Lien avec les notions du chapitre}

Cette expérience illustre :
\begin{itemize}
    \item la \textbf{polarité} : eau polaire, huile apolaire ;
    \item la \textbf{cohésion} : chaque liquide garde une cohésion propre ;
    \item la \textbf{miscibilité} : deux liquides ne sont miscibles que si les interactions croisées sont suffisamment favorables.
\end{itemize}

\begin{bilanbox}
Ce problème montre qu'un mélange homogène ne dépend pas seulement de l'action de mélanger, mais surtout des interactions microscopiques réellement possibles.
\end{bilanbox}
\end{corrigebox}

%------------------------------------------------
\subsection{Problème 2 -- Relier spectre IR et comportement d'une molécule}
%------------------------------------------------

\begin{enoncebox}
On étudie une espèce organique liquide inconnue.

On sait les choses suivantes :
\begin{itemize}
    \item son spectre IR présente une large bande vers \(3350\ \si{cm^{-1}}\) ;
    \item aucune bande intense n'apparaît vers \(1700\ \si{cm^{-1}}\) ;
    \item cette espèce est très soluble dans l'eau ;
    \item sa température d'ébullition est relativement élevée pour une molécule de petite taille.
\end{itemize}

\begin{enumerate}[label=\arabic*)]
    \item Que suggère le spectre IR ?
    \item Quelle famille de composés peut-on proposer raisonnablement ?
    \item Expliquer pourquoi les deux autres informations expérimentales sont cohérentes avec cette interprétation.
    \item Montrer que l'on relie ici directement analyse spectrale et cohésion de la matière.
\end{enumerate}
\end{enoncebox}

\begin{corrigebox}
\textbf{a) Interprétation du spectre IR}

La large bande vers \(3350\ \si{cm^{-1}}\) suggère un groupe O--H.  
L'absence de bande forte vers \(1700\ \si{cm^{-1}}\) indique l'absence probable de groupe carbonyle.

\textbf{b) Famille proposée}

On peut raisonnablement proposer un \textbf{alcool}.

\textbf{c) Cohérence avec les autres informations}

\begin{itemize}
    \item \textbf{Forte solubilité dans l'eau :} un alcool possède un groupe O--H capable d'interagir favorablement avec l'eau, notamment via des liaisons hydrogène.
    \item \textbf{Température d'ébullition relativement élevée :} les molécules d'alcool peuvent former des interactions intermoléculaires plus fortes que de simples interactions de Van der Waals, ce qui renforce la cohésion du liquide.
\end{itemize}

\textbf{d) Lien entre les deux chapitres}

Le spectre IR révèle une information structurale : présence probable d'un groupe O--H.  
Cette structure permet des liaisons hydrogène.  
Ces interactions expliquent ensuite les propriétés macroscopiques observées : solubilité élevée dans l'eau et température d'ébullition relativement importante.

\begin{ideebox}
Le spectre IR n'est pas seulement une méthode de détection de groupes chimiques ; il devient un point d'entrée pour prédire des propriétés physiques.
\end{ideebox}
\end{corrigebox}

%------------------------------------------------
\subsection{Problème 3 -- Deux molécules, deux comportements}
%------------------------------------------------

\begin{enoncebox}
On compare deux espèces :
\begin{itemize}
    \item le dioxyde de carbone \(\mathrm{CO_2}\),
    \item l'eau \(\mathrm{H_2O}\).
\end{itemize}

Toutes deux contiennent l'élément oxygène, mais leurs comportements sont très différents.

On demande :
\begin{enumerate}[label=\arabic*)]
    \item d'expliquer pourquoi leurs liaisons avec l'oxygène sont polaires ;
    \item d'expliquer pourquoi l'une des molécules est polaire et l'autre non ;
    \item d'en déduire des différences attendues dans leurs interactions intermoléculaires ;
    \item d'expliquer pourquoi ce simple changement de géométrie a de grandes conséquences physiques.
\end{enumerate}
\end{enoncebox}

\begin{corrigebox}
\textbf{a) Polarité des liaisons}

L'oxygène est plus électronégatif que l'hydrogène et que le carbone.  
Les liaisons O--H dans \(\mathrm{H_2O}\) et C=O dans \(\mathrm{CO_2}\) sont donc polaires.

\textbf{b) Polarité globale des molécules}

\begin{itemize}
    \item \(\mathrm{H_2O}\) est coudée : les effets des liaisons polaires ne se compensent pas \(\Rightarrow\) molécule polaire.
    \item \(\mathrm{CO_2}\) est linéaire : les deux effets se compensent \(\Rightarrow\) molécule apolaire.
\end{itemize}

\textbf{c) Interactions intermoléculaires attendues}

\begin{itemize}
    \item \(\mathrm{H_2O}\) : interactions fortes, notamment liaisons hydrogène.
    \item \(\mathrm{CO_2}\) : interactions plus faibles, essentiellement de Van der Waals.
\end{itemize}

\textbf{d) Pourquoi la géométrie change-t-elle tant de choses ?}

La géométrie détermine si les déséquilibres locaux de charge s'additionnent ou se compensent.  
Or la polarité globale conditionne ensuite le type d'interactions intermoléculaires possibles.  
Ces interactions gouvernent à leur tour des propriétés macroscopiques majeures.

\begin{bilanbox}
Ce problème montre que la géométrie d'une molécule n'est pas un détail géométrique : c'est une information physique structurante.
\end{bilanbox}
\end{corrigebox}

%================================================
\section{Questions plus ouvertes}
%================================================

%------------------------------------------------
\subsection{Question ouverte 1}
%------------------------------------------------

\begin{enoncebox}
Un élève dit :

\begin{center}
\textit{« En chimie, il suffit d'apprendre les tableaux de bandes IR et les règles de solubilité. »}
\end{center}

Rédiger une réponse argumentée montrant en quoi cette phrase passe à côté de l'essentiel.
\end{enoncebox}

\begin{corrigebox}
Une telle phrase réduit le chapitre à un apprentissage mécanique, alors que l'essentiel est de comprendre les liens entre structure microscopique et propriétés observables.

Apprendre un tableau de bandes IR sans comprendre le sens physique revient à reconnaître un signal sans savoir ce qu'il signifie.  
Par exemple, une bande O--H large n'est pas juste ``une case à cocher'' : elle renseigne sur la présence d'une liaison et parfois sur l'existence de liaisons hydrogène.

De même, la règle ``le semblable dissout le semblable'' n'a de valeur que si l'on comprend quelles interactions sont rompues et quelles interactions nouvelles sont créées.

La vraie compréhension consiste donc à relier :
\[
\text{structure} \longrightarrow \text{polarité} \longrightarrow \text{interactions} \longrightarrow \text{propriétés} \longrightarrow \text{signatures spectrales}
\]

Les tableaux et les règles sont des outils ; ils ne remplacent pas le raisonnement.
\end{corrigebox}

%------------------------------------------------
\subsection{Question ouverte 2}
%------------------------------------------------

\begin{enoncebox}
Expliquer, avec un exemple précis, comment une seule et même propriété macroscopique peut dépendre de plusieurs idées du cours à la fois.
\end{enoncebox}

\begin{corrigebox}
On peut prendre l'exemple de la solubilité d'un alcool dans l'eau.

Cette propriété dépend simultanément :
\begin{itemize}
    \item de la \textbf{structure moléculaire} : présence d'un groupe O--H ;
    \item de la \textbf{polarité} : ce groupe rend une partie de la molécule polaire ;
    \item des \textbf{interactions intermoléculaires} : l'alcool peut former des liaisons hydrogène avec l'eau ;
    \item de l'\textbf{analyse IR} : la présence du groupe O--H est repérable par une bande large vers \(3200-3600\ \si{cm^{-1}}\).
\end{itemize}

Ainsi, une simple observation macroscopique, ici ``cet alcool se dissout bien dans l'eau'', mobilise plusieurs étages du raisonnement physico-chimique.
\end{corrigebox}

%================================================
\section{Synthèse finale}
%================================================

\begin{bilanbox}
Dans ces deux chapitres, il faut toujours chercher à articuler cinq niveaux :

\begin{enumerate}
    \item la \textbf{nature des entités} : ions, molécules, atomes ;
    \item la \textbf{répartition des charges} : polarité locale et globale ;
    \item les \textbf{interactions possibles} : covalente, ionique, Van der Waals, liaison hydrogène ;
    \item les \textbf{conséquences macroscopiques} : solubilité, miscibilité, température d'ébullition ;
    \item les \textbf{indices expérimentaux} : bandes caractéristiques en spectroscopie IR.
\end{enumerate}

Un élève comprend vraiment le chapitre lorsqu'il est capable de passer de l'un à l'autre dans les deux sens.
\end{bilanbox}

\vfill

\begin{center}
\rule{0.82\linewidth}{0.5pt}\\
\textit{Fin de la feuille d'exercices corrigés}
\end{center}

\end{document}